\documentclass[aspectratio=169]{beamer}

\mode<presentation>
{
 \usetheme[reversetitle,notitle,noauthor]{Wien}
}

\usepackage{url}
\usepackage{graphicx}
\graphicspath{{./}{./Figures/}{./results/}}

\usepackage{appendixnumberbeamer}
\usepackage{amsmath}
\usepackage{booktabs}
\usepackage{array}
\usepackage{xcolor}

% To avoid a warning from the hyperref package:
\pdfstringdefDisableCommands{%
  \def\translate{}%
}

\definecolor{tuwred}{RGB}{140,0,0}
\definecolor{darkgreen}{RGB}{0,120,60}
\definecolor{darkblue}{RGB}{0,70,140}
\definecolor{darkorange}{RGB}{190,100,0}

\newcommand{\good}[1]{\textcolor{darkgreen}{\textbf{#1}}}
\newcommand{\warn}[1]{\textcolor{darkorange}{\textbf{#1}}}
\newcommand{\bad}[1]{\textcolor{tuwred}{\textbf{#1}}}

%%%%%%%%%%%%%%%%%%%%%%%%%%%%%%%%%%%%%%%%%%%%%%%%%%%%%%%%%%%%%%%%%%%%%%%%%%%%%
%%%%%%%%%%%%%%%%%%%%%%%%%%%%%%%%%%%%%%%%%%%%%%%%%%%%%%%%%%%%%%%%%%%%%%%%%%%%%
\title[EvoPress Attribution]{Joint Compression Attribution for EvoPress}
\subtitle{Supervisor Meeting: Progress Update and Thesis Direction}
\author[Kerim Halilovi\'c]{Kerim Halilovi\'c}
\institute[TU Wien]{TU Wien, Vienna, Austria}
\date{July 2026}

\begin{document}

\begin{frame}
  \titlepage
\end{frame}

%%%%%%%%%%%%%%%%%%%%%%%%%%%%%%%%%%%%%%%%%%%%%%%%%%%%%%%%%%%%%%%%%%%%%%%%%%%%%

\begin{frame}{Main message}
\begin{block}{Current thesis direction}
Extend EvoPress from single-method compression to \textbf{joint multi-method compression}, focusing on depth pruning + quantization.
\end{block}

\vspace{0.25cm}
\textbf{What changed since the last meeting}
\begin{itemize}
    \item The main experiments are now on \textbf{Mistral-7B}, not only TinyLlama.
    \item I implemented joint depth+quantization search and an interaction-aware mutation mode.
    \item I added a new \textbf{Joint Compression Attribution Framework}.
    \item The framework tells us whether the improvement comes from:
    \begin{itemize}
        \item the depth mask,
        \item the quantization assignment,
        \item or the exact pairing of both.
    \end{itemize}
\end{itemize}
\end{frame}

%%%%%%%%%%%%%%%%%%%%%%%%%%%%%%%%%%%%%%%%%%%%%%%%%%%%%%%%%%%%%%%%%%%%%%%%%%%%%

\begin{frame}{Thesis hypothesis}
\begin{block}{Working hypothesis}
Joint EvoPress-style search can find compression decisions that are more compatible than independently optimized pruning and quantization decisions.
\end{block}

\vspace{0.25cm}
\textbf{Why this is plausible}
\begin{itemize}
    \item Depth pruning removes complete attention or MLP modules.
    \item Quantization changes selected projection weights.
    \item These decisions interact: a good pruning mask alone may not be the best pruning mask after quantization.
    \item A joint search should be able to select candidates based on the final combined compressed model.
\end{itemize}

\vspace{0.2cm}
\begin{block}{What the attribution framework tests}
If joint search improves quality, which part of the final candidate is responsible?
\end{block}
\end{frame}

%%%%%%%%%%%%%%%%%%%%%%%%%%%%%%%%%%%%%%%%%%%%%%%%%%%%%%%%%%%%%%%%%%%%%%%%%%%%%

\begin{frame}{Why an attribution framework was needed}
\begin{columns}[T]
\begin{column}{0.48\textwidth}
\textbf{Before attribution}
\begin{itemize}
    \item We could compare final methods:
    \begin{itemize}
        \item independent compression,
        \item standard joint search,
        \item interaction-aware joint search.
    \end{itemize}
    \item But this only says which final model is better.
    \item It does not explain \textbf{why} it is better.
\end{itemize}
\end{column}

\begin{column}{0.48\textwidth}
\textbf{Question left open}
\begin{itemize}
    \item Did joint search find a better pruning mask?
    \item Did it find a better bit allocation?
    \item Did it find a pair that only works well together?
    \item Is uniform 3-bit already as good as the searched assignment?
\end{itemize}
\end{column}
\end{columns}

\vspace{0.25cm}
\begin{block}{Motivation}
The framework turns a final-result comparison into a controlled recombination analysis.
\end{block}
\end{frame}

%%%%%%%%%%%%%%%%%%%%%%%%%%%%%%%%%%%%%%%%%%%%%%%%%%%%%%%%%%%%%%%%%%%%%%%%%%%%%

\begin{frame}{What the attribution framework does}
\begin{block}{Core idea}
Take depth masks and quantization assignments from existing runs, recombine them, replay the final compressed model, and evaluate perplexity.
\end{block}

\vspace{0.2cm}
\textbf{Example recombinations}
\begin{itemize}
    \item independent depth mask + independent quant assignment
    \item joint depth mask + independent quant assignment
    \item independent depth mask + joint quant assignment
    \item joint depth mask + joint quant assignment
    \item joint depth mask + uniform 3-bit quantization
\end{itemize}

\vspace{0.2cm}
\begin{block}{Important}
This is a \textbf{post-hoc replay analysis}. It does not run a new evolutionary search and does not prove causality by itself.
\end{block}
\end{frame}

%%%%%%%%%%%%%%%%%%%%%%%%%%%%%%%%%%%%%%%%%%%%%%%%%%%%%%%%%%%%%%%%%%%%%%%%%%%%%

\begin{frame}{Attribution matrix}
\begin{center}
\begin{tabular}{c|cccc}
 & \multicolumn{4}{c}{\textbf{Quantization source}} \\
\textbf{Depth source} & independent & standard joint & interaction-aware & uniform 3-bit \\
\hline
independent & eval & eval & eval & eval \\
standard joint & eval & eval & eval & eval \\
interaction-aware & eval & eval & eval & eval \\
\end{tabular}
\end{center}

\vspace{0.3cm}
\textbf{Each cell means}
\begin{enumerate}
    \item load one depth mask,
    \item load one quantization assignment,
    \item repair/validate the active bitwidth budget,
    \item replay the combined compressed model,
    \item evaluate WikiText2, C4, and FineWeb-Edu perplexity.
\end{enumerate}

\vspace{0.15cm}
\begin{block}{Interpretation}
Rows tell us about depth-mask quality. Columns tell us about quant-assignment quality.
\end{block}
\end{frame}

%%%%%%%%%%%%%%%%%%%%%%%%%%%%%%%%%%%%%%%%%%%%%%%%%%%%%%%%%%%%%%%%%%%%%%%%%%%%%

\begin{frame}{Implementation details}
\small
\begin{itemize}
    \item Added main script:
    \begin{center}
        \texttt{evo\_joint\_attribution.py}
    \end{center}
    \item Added aggregation helper:
    \begin{center}
        \texttt{scripts/aggregate\_attribution.py}
    \end{center}
    \item Reuses existing EvoPress logic:
    \begin{itemize}
        \item candidate loading from \texttt{final\_candidate.json} or config text files,
        \item \texttt{apply\_joint\_state} for replay,
        \item \texttt{repair\_active\_quant\_budget} for active bitwidth repair,
        \item existing PPL evaluation utilities.
    \end{itemize}
    \item Outputs structured artifacts:
    \begin{itemize}
        \item combined candidate JSON,
        \item metrics JSON/CSV,
        \item compression metrics,
        \item Markdown summaries,
        \item full attribution matrix.
    \end{itemize}
\end{itemize}
\end{frame}

%%%%%%%%%%%%%%%%%%%%%%%%%%%%%%%%%%%%%%%%%%%%%%%%%%%%%%%%%%%%%%%%%%%%%%%%%%%%%

\begin{frame}{Budget handling}
\begin{block}{Problem}
When a depth mask and quantization assignment come from different runs, the active average bitwidth may no longer be exactly 3-bit.
\end{block}

\vspace{0.2cm}
\textbf{Why}
\begin{itemize}
    \item Dropped modules do not consume active quantization budget.
    \item A quant assignment can contain bitwidths for modules that are dropped by another depth mask.
    \item Recombination can therefore change the active set.
\end{itemize}

\vspace{0.2cm}
\textbf{Framework behavior}
\begin{itemize}
    \item Check active average bitwidth before replay.
    \item Repair it to target 3-bit when needed.
    \item Record whether repair happened and how many bitwidths changed.
\end{itemize}

\vspace{0.15cm}
\begin{block}{Why this matters}
All attribution cells are evaluated at matched compression, not accidentally easier or harder budgets.
\end{block}
\end{frame}

%%%%%%%%%%%%%%%%%%%%%%%%%%%%%%%%%%%%%%%%%%%%%%%%%%%%%%%%%%%%%%%%%%%%%%%%%%%%%

\begin{frame}{Experimental setup}
\begin{columns}[T]
\begin{column}{0.48\textwidth}
\textbf{Model}
\begin{itemize}
    \item \texttt{mistralai/Mistral-7B-v0.3}
    \item Depth pruning target: 25\%
    \item Quant target: active average 3-bit
    \item Datasets:
    \begin{itemize}
        \item WikiText2
        \item C4
        \item FineWeb-Edu
    \end{itemize}
\end{itemize}
\end{column}

\begin{column}{0.48\textwidth}
\textbf{Attribution cases}
\begin{itemize}
    \item \textbf{\texttt{q\_proj} scope}
    \begin{itemize}
        \item 3 seeds
        \item 12 recombinations per seed
        \item compression ratio: 1.400x
    \end{itemize}
    \item \textbf{attention scope}
    \begin{itemize}
        \item 3 seeds
        \item 12 recombinations per seed
        \item compression ratio: 1.547x
    \end{itemize}
\end{itemize}
\end{column}
\end{columns}
\end{frame}

%%%%%%%%%%%%%%%%%%%%%%%%%%%%%%%%%%%%%%%%%%%%%%%%%%%%%%%%%%%%%%%%%%%%%%%%%%%%%

\begin{frame}{Result 1: \texttt{q\_proj} attribution}
\scriptsize
\begin{table}
\centering
\setlength{\tabcolsep}{3.5pt}
\begin{tabular}{llrrr}
\toprule
Depth source & Quant source & WikiText2 & C4 & FineWeb-Edu \\
\midrule
indep. & indep. & 11.96 & 14.89 & 13.15 \\
indep. & standard joint & 11.96 & 14.94 & 13.15 \\
indep. & interaction-aware & 12.13 & 15.09 & 13.36 \\
indep. & uniform 3-bit & 11.96 & 14.88 & 13.17 \\
\midrule
standard joint & indep. & 11.21 & 14.36 & 12.43 \\
standard joint & standard joint & 11.24 & 14.39 & 12.46 \\
standard joint & interaction-aware & 11.49 & 14.61 & 12.74 \\
standard joint & uniform 3-bit & 11.21 & 14.38 & 12.44 \\
\midrule
interaction-aware & indep. & \good{10.95} & \good{13.94} & \good{12.05} \\
interaction-aware & standard joint & 10.98 & 13.97 & 12.08 \\
interaction-aware & interaction-aware & 11.09 & 14.13 & 12.22 \\
interaction-aware & uniform 3-bit & \good{10.95} & 13.97 & 12.06 \\
\bottomrule
\end{tabular}
\end{table}

\vspace{-0.05cm}
\begin{itemize}
    \item Full 3-by-4 attribution matrix; values are means over seeds 0, 1, and 2.
    \item All rows are matched at compression ratio 1.400x and active average 3-bit.
    \item Best rows all use the interaction-aware depth mask; quant source matters less.
\end{itemize}
\end{frame}

%%%%%%%%%%%%%%%%%%%%%%%%%%%%%%%%%%%%%%%%%%%%%%%%%%%%%%%%%%%%%%%%%%%%%%%%%%%%%

\begin{frame}{Interpretation: what improved?}
\begin{block}{Main attribution finding}
For \texttt{q\_proj} quantization, the \textbf{interaction-aware depth mask} is the main improvement.
\end{block}

\vspace{0.2cm}
\textbf{Evidence}
\begin{itemize}
    \item Using the interaction-aware depth mask improves mean PPL across all three datasets.
    \item It works with several quantization sources:
    \begin{itemize}
        \item independent quantization,
        \item standard joint quantization,
        \item interaction-aware quantization,
        \item uniform 3-bit.
    \end{itemize}
    \item Uniform 3-bit is very competitive with searched bit assignments.
\end{itemize}

\vspace{0.15cm}
\begin{block}{Conservative conclusion}
The framework suggests that the search mainly found a more quantization-compatible pruning mask, not a clearly superior mixed-bit assignment.
\end{block}
\end{frame}

%%%%%%%%%%%%%%%%%%%%%%%%%%%%%%%%%%%%%%%%%%%%%%%%%%%%%%%%%%%%%%%%%%%%%%%%%%%%%

\begin{frame}{Result 2: attention-scope attribution}
\scriptsize
\begin{table}
\centering
\setlength{\tabcolsep}{3.5pt}
\begin{tabular}{llrrr}
\toprule
Depth source & Quant source & WikiText2 & C4 & FineWeb-Edu \\
\midrule
indep. & indep. & 13.15 & 16.10 & 14.34 \\
indep. & standard joint & 13.20 & 16.26 & 14.38 \\
indep. & interaction-aware & 14.19 & 17.40 & 15.43 \\
indep. & uniform 3-bit & 13.13 & 16.19 & 14.29 \\
\midrule
standard joint & indep. & 12.70 & 15.63 & 13.78 \\
standard joint & standard joint & 12.67 & 15.59 & 13.78 \\
standard joint & interaction-aware & 13.72 & 17.15 & 15.02 \\
standard joint & uniform 3-bit & 12.61 & 15.60 & 13.70 \\
\midrule
interaction-aware & indep. & 12.08 & \good{14.90} & 13.02 \\
interaction-aware & standard joint & 12.30 & 15.14 & 13.12 \\
interaction-aware & interaction-aware & 12.42 & 15.25 & 13.45 \\
interaction-aware & uniform 3-bit & \good{12.07} & 14.99 & \good{12.97} \\
\bottomrule
\end{tabular}
\end{table}

\vspace{-0.05cm}
\begin{itemize}
    \item Full 3-by-4 attribution matrix; values are means over seeds 0, 1, and 2.
    \item All rows are matched at compression ratio 1.547x and active average 3-bit.
    \item Interaction-aware depth is best overall, but the interaction-aware quant assignment is weak when paired with other depth masks.
\end{itemize}
\end{frame}

%%%%%%%%%%%%%%%%%%%%%%%%%%%%%%%%%%%%%%%%%%%%%%%%%%%%%%%%%%%%%%%%%%%%%%%%%%%%%

\begin{frame}{What this supports and does not support}
\begin{columns}[T]
\begin{column}{0.48\textwidth}
\begin{block}{Supported}
\begin{itemize}
    \item Joint compression can be replayed and analyzed systematically.
    \item Interaction-aware \texttt{q\_proj} search finds a better depth mask.
    \item The improvement is consistent across WikiText2, C4, and FineWeb-Edu.
    \item Attribution helps explain the mechanism behind the result.
\end{itemize}
\end{block}
\end{column}

\begin{column}{0.48\textwidth}
\begin{block}{Not supported}
\begin{itemize}
    \item Not evidence that joint search always wins.
    \item Not evidence that searched mixed-bit assignments are always better than uniform 3-bit.
    \item Not yet evidence for full all-linear Mistral compression.
    \item Interaction-aware quant assignments are not consistently better than uniform 3-bit.
\end{itemize}
\end{block}
\end{column}
\end{columns}
\end{frame}

%%%%%%%%%%%%%%%%%%%%%%%%%%%%%%%%%%%%%%%%%%%%%%%%%%%%%%%%%%%%%%%%%%%%%%%%%%%%%

\begin{frame}{How this shapes the thesis proposal}
\begin{block}{Proposed thesis framing}
EvoPress-style joint multi-method compression for LLMs, with attribution analysis of how pruning and quantization decisions interact.
\end{block}

\vspace{0.25cm}
\textbf{Methodological contribution}
\begin{itemize}
    \item Implement joint depth+quantization search.
    \item Add interaction-aware mutation as a search-operator extension.
    \item Add replay-based attribution matrix to analyze final candidates.
\end{itemize}

\vspace{0.2cm}
\textbf{Experimental contribution}
\begin{itemize}
    \item Matched-compression Mistral-7B experiments.
    \item Multi-dataset perplexity evaluation.
    \item Evidence that the main \texttt{q\_proj} gain comes from the depth mask.
    \item Stress-test showing that broader attention quantization is harder.
\end{itemize}
\end{frame}

%%%%%%%%%%%%%%%%%%%%%%%%%%%%%%%%%%%%%%%%%%%%%%%%%%%%%%%%%%%%%%%%%%%%%%%%%%%%%

\begin{frame}{Current thesis claim}
\begin{block}{Careful claim}
Joint EvoPress-style depth+quantization search improves matched independent compression in perplexity for Mistral-7B \texttt{q\_proj} quantization. Attribution suggests that the improvement mainly comes from a more quantization-compatible depth mask. The benefit depends on quantization scope and evaluation metric.
\end{block}

\vspace{0.25cm}
\textbf{Why this is a stronger story now}
\begin{itemize}
    \item We are not only reporting final numbers.
    \item We can explain which component of the joint candidate mattered.
    \item We can separate successful cases from limitations.
    \item This gives the thesis a clearer methodology chapter:
    \begin{itemize}
        \item joint search,
        \item interaction-aware mutation,
        \item attribution replay.
    \end{itemize}
\end{itemize}
\end{frame}

%%%%%%%%%%%%%%%%%%%%%%%%%%%%%%%%%%%%%%%%%%%%%%%%%%%%%%%%%%%%%%%%%%%%%%%%%%%%%

\begin{frame}{Next steps}
\begin{enumerate}
    \item Write the algorithm section in thesis form:
    \begin{itemize}
        \item candidate representation,
        \item mutation operators,
        \item active budget repair,
        \item replay attribution matrix.
    \end{itemize}
    \item Decide whether to run attention attribution for more seeds.
    \item Add a compact downstream evaluation discussion:
    \begin{itemize}
        \item LM-eval is more mixed than perplexity,
        \item avoid overclaiming.
    \end{itemize}
    \item Turn current results into final thesis tables.
    \item Continue writing methodology and experimental setup chapters.
\end{enumerate}
\end{frame}

%%%%%%%%%%%%%%%%%%%%%%%%%%%%%%%%%%%%%%%%%%%%%%%%%%%%%%%%%%%%%%%%%%%%%%%%%%%%%

\appendix

\begin{frame}{Backup: one joint-search candidate}
\begin{block}{Candidate = depth mask + quantization assignment}
The evolutionary search stores both compression decisions in one candidate.
\end{block}

\vspace{0.2cm}
\textbf{Depth mask}
\begin{itemize}
    \item For each transformer block: keep/drop attention.
    \item For each transformer block: keep/drop MLP.
\end{itemize}

\vspace{0.15cm}
\textbf{Quantization assignment}
\begin{itemize}
    \item For selected projection modules: assign bitwidth 2, 3, or 4.
    \item Bitwidths can exist for modules later dropped by the depth mask.
    \item Active-budget computation ignores dropped modules.
\end{itemize}
\end{frame}

%%%%%%%%%%%%%%%%%%%%%%%%%%%%%%%%%%%%%%%%%%%%%%%%%%%%%%%%%%%%%%%%%%%%%%%%%%%%%

\begin{frame}{Backup: one generation of joint search}
\small
\begin{enumerate}
    \item Start from the current parent candidate.
    \item Copy it into offspring candidates.
    \item Mutate candidates:
    \begin{itemize}
        \item standard mode: mutate depth or quantization,
        \item interaction-aware mode: mutate depth, repair active budget, then adjust quantization on touched active layers.
    \end{itemize}
    \item Replay each compressed candidate.
    \item Evaluate KL divergence against dense model logits.
    \item Select the best candidate as the next parent.
\end{enumerate}

\vspace{0.15cm}
\begin{block}{Key distinction}
Implementation may apply quantization before pruning for efficiency, but selection is based on the final combined compressed model.
\end{block}
\end{frame}

%%%%%%%%%%%%%%%%%%%%%%%%%%%%%%%%%%%%%%%%%%%%%%%%%%%%%%%%%%%%%%%%%%%%%%%%%%%%%

\begin{frame}{Backup: interaction-aware mutation}
\small
\begin{enumerate}
    \item Copy the current candidate.
    \item Make a budget-preserving depth swap.
    \item Recompute the active quantized modules under the new mask.
    \item Repair active average bitwidth to the target.
    \item Prefer quantization changes on layers touched by the depth swap.
    \item Log:
    \begin{itemize}
        \item mutation mode,
        \item depth changes,
        \item quantization changes,
        \item repair count.
    \end{itemize}
\end{enumerate}

\vspace{0.15cm}
\begin{block}{Implementation contribution}
This changes the search operator, not the model architecture, dataset, or quantization database.
\end{block}
\end{frame}

%%%%%%%%%%%%%%%%%%%%%%%%%%%%%%%%%%%%%%%%%%%%%%%%%%%%%%%%%%%%%%%%%%%%%%%%%%%%%

\begin{frame}{Backup: final-result comparison before attribution}
\small
\begin{table}
\centering
\begin{tabular}{llrrrr}
\toprule
Scope & Method & Compression & WikiText2 & C4 & FineWeb-Edu \\
\midrule
baseline & Dense FP16 & 1.000x & \good{5.960} & \good{8.860} & \good{7.340} \\
depth & Depth-only 25\% & 1.317x & 11.817 & 14.760 & 12.980 \\
\texttt{q\_proj} & Independent depth + quant & 1.400x & 11.947 & 14.910 & 13.140 \\
\texttt{q\_proj} & Standard joint & 1.400x & 11.243 & 14.393 & 12.460 \\
\texttt{q\_proj} & Interaction-aware joint & 1.400x & \good{11.087} & \good{14.133} & \good{12.223} \\
attention & Independent depth + quant & 1.547x & 13.107 & 16.090 & 14.237 \\
attention & Standard joint & 1.547x & \good{12.670} & \good{15.593} & \good{13.783} \\
\bottomrule
\end{tabular}
\end{table}

\vspace{0.1cm}
\textbf{Limitation:} this table shows final method quality, but not which part of the joint candidate caused the improvement.
\end{frame}

%%%%%%%%%%%%%%%%%%%%%%%%%%%%%%%%%%%%%%%%%%%%%%%%%%%%%%%%%%%%%%%%%%%%%%%%%%%%%

\begin{frame}{Backup: downstream task caveat}
\small
\begin{table}
\centering
\begin{tabular}{lrrrr}
\toprule
Method & ARC-Easy & PIQA & Winogrande & Macro \\
\midrule
Dense FP16 & \good{0.801} & \good{0.819} & \good{0.742} & \good{0.787} \\
Depth-only 25\% & 0.588 & 0.737 & 0.618 & 0.648 \\
\texttt{q\_proj} independent & 0.583 & 0.739 & 0.620 & 0.647 \\
\texttt{q\_proj} standard joint & 0.580 & 0.739 & 0.608 & 0.642 \\
\texttt{q\_proj} interaction-aware & 0.585 & 0.742 & 0.603 & 0.643 \\
Attention independent & 0.535 & 0.696 & 0.558 & 0.596 \\
Attention joint & 0.534 & 0.707 & 0.561 & 0.601 \\
\bottomrule
\end{tabular}
\end{table}

\vspace{0.1cm}
\textbf{Takeaway:} perplexity results are clearer than downstream task results, so the thesis should not overclaim downstream gains.
\end{frame}

%%%%%%%%%%%%%%%%%%%%%%%%%%%%%%%%%%%%%%%%%%%%%%%%%%%%%%%%%%%%%%%%%%%%%%%%%%%%%

\begin{frame}{Backup: reproducibility artifacts}
\small
\begin{itemize}
    \item Main attribution implementation:
    \begin{itemize}
        \item \texttt{evo\_joint\_attribution.py}
        \item \texttt{scripts/aggregate\_attribution.py}
        \item \texttt{results/attribution\_framework\_report.md}
    \end{itemize}
    \item Main result artifacts:
    \begin{itemize}
        \item \texttt{results/attribution/mistral\_qproj\_aggregate.md}
        \item \texttt{results/attribution/mistral\_qproj\_seed0/}
        \item \texttt{results/attribution/mistral\_qproj\_seed1/}
        \item \texttt{results/attribution/mistral\_qproj\_seed2/}
        \item \texttt{results/attribution/mistral\_attention\_seed0/}
    \end{itemize}
    \item Search-run artifacts:
    \begin{itemize}
        \item \texttt{final\_candidate.json}
        \item \texttt{run\_summary.json}
        \item \texttt{generation\_log.csv}
        \item \texttt{experiment\_log.csv}
    \end{itemize}
\end{itemize}
\end{frame}

%%%%%%%%%%%%%%%%%%%%%%%%%%%%%%%%%%%%%%%%%%%%%%%%%%%%%%%%%%%%%%%%%%%%%%%%%%%%%

\begin{frame}{Backup: hardware context}
\begin{itemize}
    \item Datalab GPU varied depending on availability:
    \begin{itemize}
        \item T4,
        \item V100,
        \item A40.
    \end{itemize}
    \item The main bottleneck was often the 16 GB container CPU RAM cap.
    \item Full Mistral sparse-database generation remained impractical in this environment.
    \item Restricted Mistral quantization scopes were feasible:
    \begin{itemize}
        \item \texttt{q\_proj},
        \item attention projections.
    \end{itemize}
\end{itemize}
\end{frame}

\end{document}
